\chapter{Estimación y SLAM}

El estimador alterna predicción y corrección. Primero propaga
la pose con odometría diferencial; después la corrige con la información de los
sensores (actualización EKF de rango y rumbo en \texttt{KALMAN\_LANDMARK}, o
ajuste \textit{scan-to-map} en las variantes de grid) y lleva las lecturas de
sensor a coordenadas del warehouse. Una referencia global débil y ruidosa (tipo
baliza UWB) actúa además como ancla de baja tasa para mantener observable el
marco global. Con esas detecciones se actualiza un
mapa compacto de landmarks o una grid de ocupación, según el algoritmo
seleccionado \cite{thrun2005}.

La pose que entrega CoppeliaSim tiene seis componentes:
\([x,y,z,\alpha,\beta,\gamma]\). El estimador usa solo la proyección plana
\([x,y,\theta]\), con \(\theta=\gamma\). La altura, roll y pitch se descartan
porque el Pioneer permanece sobre el piso de la celda y la tarea se formula en
2D.

\section{Sensores: LiDAR, ultrasonidos y rayos de proximidad}

Un LiDAR 2D real emite muchos pulsos láser en ángulos conocidos. Para cada ángulo
recibe una distancia, normalmente por tiempo de vuelo, y forma un escaneo polar:
\[
\mathcal{S}_k=\{(r_i,\beta_i)\}_{i=1}^{N}.
\]
Cada par \((r_i,\beta_i)\) da la distancia al primer obstáculo en ese ángulo.
Con la pose del robot, ese punto se transforma al marco global:
\[
p_i^W =
\begin{bmatrix}\hat{x}\\\hat{y}\end{bmatrix}+
R(\hat{\theta})
\begin{bmatrix}
r_i\cos\beta_i\\
r_i\sin\beta_i
\end{bmatrix}.
\]

El Pioneer usado aquí no lleva un LiDAR físico 2D. La escena usa los 16 sensores
ultrasónicos/proximidad del modelo P3DX, dibujados como rayos para que se vea
qué detecta cada sensor. Conceptualmente se tratan igual que un escaneo de baja
resolución: cada sensor tiene un ángulo fijo \(\beta_i\), devuelve una distancia
\(r_i\) si hay impacto y se transforma al mapa con la pose estimada. La
diferencia está en la densidad: un LiDAR aporta muchos más rayos y contornos
más continuos. El arreglo del Pioneer tiene pocas direcciones, de modo que los
mapas son aproximaciones locales. Las variantes GMapping, Hector y Cartographer
se implementan sobre esos rayos de proximidad simulados para comparar
representaciones de mapa, sin equivalencia con sus paquetes ROS completos.
La Fig.~\ref{fig:estimator-pipeline} muestra el ciclo completo desde sensores hasta control.

\begin{figure}[H]
  \centering
  \includegraphics[width=\textwidth]{phase1_estimator_pipeline.png}
  \caption{Secuencia de estimación, mapa, planificación local y control de R1.}
  \label{fig:estimator-pipeline}
\end{figure}

\section{Datos usados para comparar SLAM}

Los cuatro métodos reciben las mismas entradas: odometría diferencial, pose
ruidosa de referencia para cerrar el estimador, rayos de proximidad y el mismo
objetivo de navegación. La comparación evalúa cuatro representaciones implementadas en Lua para la misma
celda: rasgos compactos, grid de ocupación, ajuste local por scan matching y
submapas. El alcance difiere del de los paquetes ROS completos.

Las cuatro variantes reciben las mismas entradas y el mismo ruido: idéntica
odometría diferencial, la misma referencia global débil con ruido gaussiano
(\(\sigma=\SI{0.030}{m}\) en \(x,y\) y \(\SI{0.018}{rad}\) en yaw, con ganancias
\(k_{xy}=0{,}14\) y \(k_\theta=0{,}20\)) y los mismos rayos de proximidad. Lo que
cambia es cómo usa cada método la información del sensor para corregir la pose y
construir el mapa:

\begin{itemize}
  \item \texttt{KALMAN\_LANDMARK}: actualización EKF de la pose por rango y rumbo
  a cada landmark asociado, con gating de Mahalanobis (umbral \(\chi^2_2=9{,}0\))
  para descartar asociaciones erróneas y peso de corrección \(0{,}35\).
  \item \texttt{GMAPPING}: solo grid de ocupación; la pose proviene de odometría
  más el ancla global.
  \item \texttt{HECTOR} y \texttt{CARTOGRAPHER}: corrección de pose por ajuste
  \textit{scan-to-map} Gauss-Newton sobre la grid log-odds (ganancias de mezcla
  \(0{,}65\) y \(0{,}55\)).
\end{itemize}

La comparación es por tanto equitativa en ruido: las diferencias de
RMSE\footnote{El RMSE se mide contra la referencia global ruidosa (pose del
simulador más ruido gaussiano), no contra una verdad de terreno; es una medida de
consistencia del estimador bajo condiciones reproducibles, no de localización
absoluta.} reflejan la representación de mapa y el método de corrección, no un
ruido inyectado distinto por método.

La covarianza de medición de la actualización EKF se parametriza con
\(\sigma_r=\SI{0.075}{m}\) y \(\sigma_{\beta}=\SI{0.10}{rad}\) para rango y
rumbo. Esto mantiene repetible la comparación; no pretende modelar el ruido
exacto de un sensor ultrasónico real.

Cada algoritmo guarda algo distinto y eso cambia lo que recibe el planificador:
Kalman-landmark, pocos puntos fusionados; GMapping, probabilidades por celda;
Hector, la corrección de pose por ajuste scan-grid; Cartographer, submapas con
cierres de ciclo. Como un landmark y una celda ocupada no significan lo mismo, la
comparación no se reduce a contar rasgos.

\section{Por qué hace falta Kalman}

La odometría diferencial predice bien durante pocos instantes, pero acumula error:
un pequeño desajuste de rueda, saturación o giro se convierte en deriva de
\(x,y,\theta\). Las mediciones de sensor tampoco son perfectas; un impacto puede
aparecer desplazado por ruido, discretización o por estar mirando un borde. El
filtro de Kalman se usa para combinar esas dos fuentes en vez de creer ciegamente
en una sola.

Primero se predice la pose con el modelo de movimiento y los comandos de rueda;
luego, ante una medición ruidosa, se calcula la corrección. La ganancia \(K\)
pesa la medición según la incertidumbre: con \(P\) alta corrige más, y con \(R\)
poco fiable corrige menos. En la escena esto mantiene una pose suave para el
controlador y evita que un rayo aislado mueva el mapa de forma brusca.

En esta actividad Kalman es la opción base porque encaja con una celda pequeña
y pocos obstáculos: mantiene una pose
estable, fusiona detecciones cercanas como landmarks y entrega al planificador una lista
compacta de puntos peligrosos. Su límite: no reconstruye paredes ni superficies
completas como una grid de ocupación.

La corrección combina dos fuentes y la del sensor domina. La principal es la
observación de los sensores de proximidad: cada vez que una lectura se asocia a
un landmark del mapa, una actualización EKF corrige la pose por rango y rumbo.
El jacobiano de observación \(H\) es la derivada de la función que, dada la pose,
predice el rango y el rumbo al landmark; por eso es \(2\times3\) (dos medidas
frente a las tres componentes de la pose). La covarianza de innovación
\(S=HPH^\top+R\) es la incertidumbre esperada de esa predicción (la de la pose
llevada al sensor, más el ruido \(R\)); dividir por \(S\) en \(K=PH^\top S^{-1}\)
hace que una innovación poco fiable pese menos. Antes de aplicar la corrección,
un \emph{gate} de Mahalanobis filtra las asociaciones malas: esa distancia mide
cuánto se aparta la innovación pesándola con \(S^{-1}\) y, al tener dos
componentes (rango y rumbo), sigue una \(\chi^2\) con 2 grados de libertad; el
umbral \(9{,}0\) equivale a un nivel de confianza cercano al 98\,\%, así que por
encima la asociación se descarta. Esto evita el fallo típico que hace diverger un
EKF-SLAM con sónares dispersos. El peso de corrección aplicado es \(0{,}35\). La segunda fuente es una referencia global débil con ruido gaussiano
genuino, aplicada con ganancia pequeña solo para mantener observable el marco
global (en especial el yaw) en misiones largas. Trabajar con la covarianza
\(3\times3\) de la pose, en vez de una covarianza conjunta pose-mapa, es lo que
evita la divergencia con un anillo de 16 ultrasonidos; coincide con el alcance de
\texttt{KALMAN\_LANDMARK}, que no mantiene la covarianza conjunta de un EKF-SLAM
completo \cite{thrun2005}.

\section{Predicción de pose}

Con velocidades odométricas \(v_k,\omega_k\), el modelo diferencial discreto es:

\[
\begin{aligned}
\hat{x}_{k}^{-} &= \hat{x}_{k-1}^{+} + v_k\cos(\hat{\theta}_{k-1}^{+})\Delta t\\
\hat{y}_{k}^{-} &= \hat{y}_{k-1}^{+} + v_k\sin(\hat{\theta}_{k-1}^{+})\Delta t\\
\hat{\theta}_{k}^{-} &= \operatorname{wrap}(\hat{\theta}_{k-1}^{+}+\omega_k\Delta t)
\end{aligned}
\]

La incertidumbre también se propaga. Para ello se deriva el modelo de movimiento
respecto a la pose anterior: ese jacobiano \(G_k\) tiene como únicos términos no
triviales \(-v_k\Delta t\sin\theta\) y \(v_k\Delta t\cos\theta\), que miden cuánto
se desplazan \(x\) e \(y\) ante un pequeño error de orientación. La covarianza
predicha es entonces \(P_k^-=G_k P_{k-1}^+ G_k^\top+Q_k\) (la del paso anterior
arrastrada por \(G_k\), más el ruido de proceso \(Q_k\)):

\[
G_k =
\begin{bmatrix}
1 & 0 & -v_k\Delta t\sin\hat{\theta}_{k-1}^{+}\\
0 & 1 &  v_k\Delta t\cos\hat{\theta}_{k-1}^{+}\\
0 & 0 & 1
\end{bmatrix},
\qquad
P_k^- = G_k P_{k-1}^+ G_k^\top + Q_k
\]

\section{Corrección Kalman-landmark}

La formulación de referencia sigue el filtro de Kalman extendido (EKF) para sistemas de navegación con landmarks \cite{thrun2005}:

\[
K_k = P_k^-H_k^\top(H_kP_k^-H_k^\top + R_k)^{-1}
\]

\[
\hat{x}_k^+ = \hat{x}_k^- + K_k(z_k-H_k\hat{x}_k^-),\qquad
P_k^+ = (I-K_kH_k)P_k^-
\]

En la escena, esta actualización corrige la pose del robot a partir de cada
landmark observado (rango y rumbo), con el \emph{gate} de Mahalanobis que rechaza
asociaciones improbables. Las detecciones se asocian por cercanía: si una lectura
cae dentro del radio de asociación de un landmark, su posición se refina y la
pose se corrige; las lecturas no asociadas crean landmarks nuevos, hasta el
máximo configurado. La posición de un landmark nuevo se obtiene invirtiendo la
observación: a partir de la pose y del par (rango, rumbo) se calcula su punto
\((x,y)\). Como esa posición hereda dos incertidumbres (la de la pose y la del
sensor), se propaga cada una con la derivada de esa función inversa: \(G_R\)
respecto a la pose y \(G_z\) respecto a la medición. De ahí los dos sumandos de
\(P_{LL}=G_R P_{RR} G_R^\top+G_z R G_z^\top\), más la covarianza cruzada con la
pose. \texttt{KALMAN\_LANDMARK} se denomina así
porque corrige la pose con la covarianza \(3\times3\) del robot frente al mapa de
rasgos, sin mantener la covarianza conjunta completa pose-mapa de un EKF-SLAM
total (Fig.~\ref{fig:slam-kalman-diagram}).

\begin{figure}[H]
\centering
\begin{tikzpicture}[
  >=Latex,
  node distance=0.70cm and 0.65cm,
  block/.style={draw=UniPrimaryDeep!80, fill=UniLight, rounded corners=2mm, minimum width=3.0cm, minimum height=0.82cm, align=center, font=\small},
  map/.style={draw=GoodGreen!80!black, fill=GoodGreen!8, rounded corners=2mm, minimum width=3.0cm, minimum height=0.82cm, align=center, font=\small},
  note/.style={draw=WarnAmber!85!black, fill=WarnAmber!10, rounded corners=2mm, minimum width=3.0cm, minimum height=0.72cm, align=center, font=\small}
]
\node[block] (odom) {Odometría\\\(v_k,\omega_k,\Delta t\)};
\node[block, right=of odom] (pred) {Predicción\\\(\hat{x}_k^-,P_k^-\)};
\node[block, right=of pred] (innov) {Innovación\\\(z_k-H\hat{x}_k^-\)};
\node[block, right=of innov] (corr) {Corrección\\\(K_k\), \(\hat{x}_k^+\)};
\node[map, below=of innov] (lm) {Mapa compacto\\landmarks fusionadas};
\node[note, below=of pred] (assoc) {Asociación\\por cercanía};
\draw[->, thick, color=UniPrimaryDeep] (odom) -- (pred);
\draw[->, thick, color=UniPrimaryDeep] (pred) -- (innov);
\draw[->, thick, color=UniPrimaryDeep] (innov) -- (corr);
\draw[->, thick, color=GoodGreen!70!black] (corr) |- (lm);
\draw[->, thick, color=WarnAmber!80!black] (lm) -- (assoc);
\draw[->, thick, color=WarnAmber!80!black] (assoc) -- (innov);
\end{tikzpicture}
\caption{Particularidad de Kalman-landmark: corrige la pose y fusiona pocas detecciones como rasgos; no mantiene una covarianza conjunta pose-mapa como un EKF-SLAM completo.}
\label{fig:slam-kalman-diagram}
\end{figure}

En un entorno pequeño, esta opción mantiene estabilidad: al fusionar
lecturas cercanas, el mapa no crece sin control y el planificador recibe
obstáculos puntuales fáciles de filtrar. La desventaja es visual: no reconstruye
superficies completas, de modo que una pared larga aparece como pocos rasgos y
no como una frontera continua.

\section{GMapping en grid}

GMapping se basa en filtros de partículas Rao-Blackwellizados y mapas de grid
\cite{grisetti2007}. En esta simulación se implementó una grid de ocupación
con log-odds, donde \(L(c)=\log(p(c)/(1-p(c)))\). Para cada rayo, las celdas
atravesadas se marcan como libres y la celda de impacto como ocupada:

\[
L_t(c)=\operatorname{clip}\left(L_{t-1}(c)+
\begin{cases}
l_{occ}, & c=c_{hit}\\
l_{free}, & c\in ray
\end{cases}\right)
\]

Los incrementos son \(l_{occ}=0{,}85\) y \(l_{free}=-0{,}18\); se acotan con
\(\operatorname{clip}(\cdot, -3{,}60,\, 3{,}60)\) para mantener las probabilidades
en el rango \([0{,}027,\, 0{,}973]\) y evitar saturación por lecturas aisladas.
Una celda se considera ocupada si \(L(c)\geq 1{,}05\) (parámetro \texttt{gridOccupiedThreshold}).

La grid resultante se usa tanto para telemetría como para el planificador
local. En Lua, el recorrido de celdas se obtiene muestreando el
segmento entre el sensor y el punto detectado con paso menor que la resolución
de la grid. Es una variante simple de ray-casting, suficiente para marcar las
celdas libres antes de la celda de impacto (Fig.~\ref{fig:slam-gmapping-diagram}).

\begin{figure}[H]
\centering
\begin{tikzpicture}[
  >=Latex,
  node distance=0.72cm and 0.75cm,
  block/.style={draw=UniPrimaryDeep!80, fill=UniLight, rounded corners=2mm, minimum width=2.85cm, minimum height=0.78cm, align=center, font=\small},
  grid/.style={draw=GoodGreen!80!black, fill=GoodGreen!8, rounded corners=2mm, minimum width=2.85cm, minimum height=0.78cm, align=center, font=\small},
  warn/.style={draw=WarnAmber!85!black, fill=WarnAmber!10, rounded corners=2mm, minimum width=2.85cm, minimum height=0.78cm, align=center, font=\small}
]
\node[block] (scan) {Rayos de sensor\\alcance + impacto};
\node[block, right=of scan] (ray) {Ray casting\\celdas atravesadas};
\node[grid, right=of ray] (logodds) {Log-odds\\libre / ocupada};
\node[grid, right=of logodds] (map) {Grid SLAM\\celdas activas};
\node[warn, below=of logodds] (plan) {Planificador local\\bloqueos por celda};
\draw[->, thick, color=UniPrimaryDeep] (scan) -- (ray);
\draw[->, thick, color=UniPrimaryDeep] (ray) -- (logodds);
\draw[->, thick, color=GoodGreen!70!black] (logodds) -- (map);
\draw[->, thick, color=WarnAmber!80!black] (map) -- (plan);
\draw[->, thick, color=WarnAmber!80!black] (plan.west) -| (scan.south);
\end{tikzpicture}
\caption{Particularidad de GMapping implementado: la información principal es una grid de ocupación. Cada rayo libera celdas antes del impacto y refuerza la celda ocupada.}
\label{fig:slam-gmapping-diagram}
\end{figure}

El GMapping original usa un filtro de partículas Rao-Blackwellizado para
mantener hipótesis de trayectoria. En esta actividad se conservó la idea que
interesa para el planificador, la grid de ocupación, y se redujo el resto para
que el coste computacional fuera razonable dentro del script de CoppeliaSim. Esa
regla explica por qué GMapping genera muchos más elementos de mapa que
Kalman-landmark, pero no necesariamente menor error de pose.

\section{Hector con ajuste local}

Hector SLAM usa \textit{scan matching} contra una grid y depende menos de la
odometría \cite{kohlbrecher2011}. La implementación resuelve el ajuste
scan-mapa mediante el método de Gauss-Newton con jacobianos analíticos,
en lugar de una búsqueda exhaustiva.

La función de puntuación usa interpolación bilineal del mapa log-odds para
obtener valores y gradientes continuos:

\[
M_{bilin}(x,y)=\sum_{k\in\{00,10,01,11\}} \lambda_k(x,y)\,L(c_k),
\quad
\nabla M_{bilin}=\frac{1}{\Delta}\begin{bmatrix}\partial_x\\\partial_y\end{bmatrix}M_{bilin}
\]

donde \(\Delta\) es la resolución de la grid, \(\lambda_k\) son los pesos
bilineales y \(L(c_k)\) el valor log-odds de la celda vecina \(c_k\).
La pose óptima maximiza la puntuación media del scan:
\[
\xi^\star=\arg\max_\xi \frac{1}{N}\sum_{i=1}^{N}M_{bilin}\!\left(T_\xi(p_i)\right).
\]

Para cada punto del scan se calcula el jacobiano analítico de la
transformación rígida \(T_\xi(p_i)=(x+r_i\cos(\theta+\beta_i),\;y+r_i\sin(\theta+\beta_i))\):

\[
J_i=\nabla M_{bilin}\!\begin{bmatrix}1 & 0 & -r_i\sin(\theta+\beta_i)\\0 & 1 & r_i\cos(\theta+\beta_i)\end{bmatrix}
\in\mathbb{R}^{1\times3}
\]

Gauss-Newton aproxima el Hessiano como \(H\mathrel{+}=J_i^TJ_i\): desprecia las
segundas derivadas del mapa y se queda solo con los gradientes \(J_i\) que ya
calculamos, lo que basta cerca del óptimo. El paso \(\Delta\xi\) sale de resolver
\(H\,\Delta\xi=b\) (con \(b=\tfrac{1}{N}\sum_i J_i^TM_i\)); como el sistema es
apenas \(3\times3\), se resuelve con la regla de Cramer (determinantes), sin
factorizaciones. Se aplican hasta 5 iteraciones, con convergencia
\(\|\Delta\xi\|<0{,}8\,\text{mm}\) y \(|\Delta\theta|<1{,}5\,\text{mrad}\).
Esta formulación sigue directamente la derivación de Hector SLAM
\cite{kohlbrecher2011}, sin las múltiples resoluciones del paquete ROS completo
(Fig.~\ref{fig:slam-hector-diagram}).

\begin{figure}[H]
\centering
\begin{tikzpicture}[
  >=Latex,
  node distance=0.68cm and 0.62cm,
  block/.style={draw=UniPrimaryDeep!80, fill=UniLight, rounded corners=2mm, minimum width=2.90cm, minimum height=0.78cm, align=center, font=\small},
  score/.style={draw=GoodGreen!80!black, fill=GoodGreen!8, rounded corners=2mm, minimum width=2.90cm, minimum height=0.78cm, align=center, font=\small},
  warn/.style={draw=WarnAmber!85!black, fill=WarnAmber!10, rounded corners=2mm, minimum width=2.90cm, minimum height=0.78cm, align=center, font=\small}
]
\node[block] (pose) {Pose \(\hat\xi=(\hat x,\hat y,\hat\theta)\)};
\node[block, right=of pose] (interp) {Interpolar\\\(M_{bilin},\nabla M_{bilin}\)};
\node[score, right=of interp] (hess) {Acumular\\\(H\mathrel{+}=J_i^TJ_i\)};
\node[warn, below=of hess] (solve) {Resolver\\\(H\Delta\xi=b\)};
\node[warn, left=of solve] (update) {Actualizar\\\(\hat\xi\mathrel{+}=\Delta\xi\)};
\draw[->, thick, color=UniPrimaryDeep] (pose) -- (interp);
\draw[->, thick, color=UniPrimaryDeep] (interp) -- (hess);
\draw[->, thick, color=GoodGreen!70!black] (hess) -- (solve);
\draw[->, thick, color=WarnAmber!80!black] (solve) -- (update);
\draw[->, thick, color=WarnAmber!80!black] (update) -- (pose);
\end{tikzpicture}
\caption{Hector SLAM: iteración Gauss-Newton con interpolación bilineal del mapa log-odds. El Hessiano se acumula con los jacobianos analíticos de la transformación rígida; el sistema \(H\Delta\xi=b\) se resuelve por la regla de Cramer.}
\label{fig:slam-hector-diagram}
\end{figure}

La interpolación bilineal da gradientes continuos para el Gauss-Newton, frente a
una evaluación discreta de celdas. Si el entorno
tiene poca textura (pasillos simétricos), el Hessiano puede ser casi singular;
la regularización diagonal \(0{,}002\) previene divergencia y el umbral de
puntuación mínima (\texttt{hectorMatchMinScore}) descarta correcciones débiles.

\section{Cartographer con submapas}

Google Cartographer organiza el mapeo en submapas y cierres de ciclo para
reducir deriva acumulada \cite{hess2016}. La implementación crea submapas
por distancia recorrida (\(\SI{1.85}{m}\)) o periodo de tiempo (\(\SI{18}{s}\)),
detecta cierres de ciclo cuando el robot regresa cerca de un submapa antiguo
y optimiza el grafo de poses para propagar la corrección globalmente.

Un submapa se considera candidato a cierre si lleva más de \(\SI{22}{s}\)
creado, la distancia es inferior a \(\SI{0.45}{m}\), la separación en el
grafo es de al menos tres submapas y el puntaje de scan-matching supera
\(0{,}35\). Cuando se detecta un cierre válido, se registra la
restricción como arista del grafo de poses:

\[
e_{ij} = \bigl(\Delta x_{ij},\,\Delta y_{ij},\,\Delta\theta_{ij}\bigr)
\]

donde $(\Delta x_{ij},\Delta y_{ij},\Delta\theta_{ij})$ es la pose relativa
medida entre el submapa activo $j$ y el submapa antiguo $i$ en el instante del
cierre. La optimización del grafo minimiza el error cuadrático total sobre
todas las restricciones almacenadas mediante descenso Gauss-Seidel:

\[
\varepsilon_{ij}=\bigl((\hat x_j-\hat x_i)-\Delta x_{ij},\;
(\hat y_j-\hat y_i)-\Delta y_{ij},\;
\operatorname{wrap}((\hat\theta_j-\hat\theta_i)-\Delta\theta_{ij})\bigr)
\]

\[
\hat x_i \mathrel{+}=\lambda\,\varepsilon_{ij,x},\quad
\hat x_j \mathrel{-}=\lambda\,\varepsilon_{ij,x},\quad \lambda=0{,}14,\quad\text{(ídem }y,\theta\text{)}
\]

El residuo \(\varepsilon_{ij}\) es cuánto difiere la pose relativa estimada entre
los dos submapas de la que mide el cierre. Ese error se reparte moviendo ambas
poses una fracción \(\lambda=0{,}14\) en sentidos opuestos (una se acerca y la
otra se aleja), de modo que cada paso corrige sin saltos bruscos; Gauss-Seidel es
simplemente recorrer las restricciones de una en una reutilizando los valores ya
actualizados. Se ejecutan 8 iteraciones por cierre detectado y la pose del robot
se ancla suavemente al submapa activo tras la optimización.
En la comparación de Fase 1, Cartographer generó 14 submapas y 3 cierres de ciclo.
Esta formulación captura la idea central de Cartographer \cite{hess2016} sin
el solucionador Ceres ni la optimización multi-resolución del paquete original.

El flujo de submapas y cierres de ciclo se ilustra en la Fig.~\ref{fig:slam-cartographer-diagram}.

\begin{figure}[H]
\centering
\begin{tikzpicture}[
  >=Latex,
  node distance=0.70cm and 0.70cm,
  block/.style={draw=UniPrimaryDeep!80, fill=UniLight, rounded corners=2mm, minimum width=2.75cm, minimum height=0.78cm, align=center, font=\small},
  submap/.style={draw=GoodGreen!80!black, fill=GoodGreen!8, rounded corners=2mm, minimum width=2.75cm, minimum height=0.78cm, align=center, font=\small},
  loop/.style={draw=WarnAmber!85!black, fill=WarnAmber!10, rounded corners=2mm, minimum width=2.75cm, minimum height=0.78cm, align=center, font=\small}
]
\node[block] (scan) {Scan actual\\pose local};
\node[submap, right=of scan] (active) {Submapa activo\\acumula celdas};
\node[submap, right=of active] (new) {Submapa adicional\\por distancia/tiempo};
\node[loop, below=of active] (candidate) {Candidato antiguo\\cercanía + puntaje};
\node[loop, right=of candidate] (closure) {Cierre de ciclo\\ajuste local};
\draw[->, thick, color=UniPrimaryDeep] (scan) -- (active);
\draw[->, thick, color=GoodGreen!70!black] (active) -- (new);
\draw[->, thick, color=WarnAmber!80!black] (new) |- (candidate);
\draw[->, thick, color=WarnAmber!80!black] (candidate) -- (closure);
\draw[->, thick, color=WarnAmber!80!black] (closure) -| (active);
\end{tikzpicture}
\caption{Particularidad de Cartographer-submap: el mapa no se guarda como una sola grid monolítica, sino como submapas locales con posibles cierres cuando R1 vuelve a zonas conocidas.}
\label{fig:slam-cartographer-diagram}
\end{figure}

El mapa guarda estructura temporal: qué zona se construyó en cada tramo y los
retornos a regiones ya visitadas. El coste es una lógica de navegación local más
delicada: si una celda antigua queda cerca de la ruta, el planificador la trata
como bloqueo; por eso en Fase 2 se añadieron filtros de persistencia, distancia
al robot y recuperación directa.

\section{Qué método usar y por qué}

La Tabla~\ref{tab:slam-selection} resume cuándo conviene cada método: pose y
puntos cercanos a evitar (Kalman-landmark), ocupación libre/ocupada por celda
(GMapping), corrección de pose por encaje del escaneo en el mapa (Hector) y
submapas con reconocimiento de zonas ya visitadas (Cartographer).

\begin{table}[H]
\centering
\small
\setlength{\tabcolsep}{3pt}
\begin{tabularx}{\textwidth}{@{}P{0.20\textwidth}P{0.28\textwidth}P{0.25\textwidth}Y@{}}
\toprule
\textbf{Método} & \textbf{Conviene cuando} & \textbf{No conviene cuando} & \textbf{Lectura en esta actividad}\\
\midrule
Kalman-landmark & Se necesita pose estable y pocos obstáculos puntuales. & Se requiere una pared o contorno continuo del entorno. & Menor duración y mapa compacto; opción base para la celda compacta.\\
GMapping/log-odds & El planificador necesita saber qué celdas están ocupadas. & Se espera estimación completa de trayectoria por partículas. & Útil como grid de ocupación ligera; no implementa el GMapping completo.\\
Hector/scan-to-map & Hay suficiente geometría para alinear el scan contra el mapa. & Los rayos son pocos, simétricos o con impactos pobres. & Da buena precisión de rasgos en Fase 2, pero puede oscilar si el scan no discrimina candidatos.\\
Cartographer & El robot recorre zonas amplias y vuelve a regiones ya vistas. & La escena es pequeña y el coste de gestionar submapas supera el beneficio. & Logra mayor cobertura y cierres, pero tarda más y exige filtros para no bloquear la ruta.\\
\bottomrule
\end{tabularx}
\caption{Criterio de selección entre las variantes SLAM implementadas.}
\label{tab:slam-selection}
\end{table}

Con los datos de la actividad, la elección operativa es Kalman-landmark para
controlar el Pioneer en esta celda. Las grids y submapas se mantienen para
comparar cómo cambia el mapa cuando se prioriza ocupación, ajuste local o
estructura temporal.
